% Conclusion (main paper). Numbers: tab1_controlled, tab4_coverage, tab5_reasoning (test splits).
\section{Conclusion}
\label{sec:conclusion}

We presented \ours{}, which grounds a part by asking a multimodal LLM to choose a mask instead of drawing one. A
frozen SAM\,3 proposes candidates with and without the query, every candidate enters the model as one mask token
that reads only its own region, and a selection token scores them all in a single pass. With $64.2$M trainable
parameters and a $4$B backbone, \ours{} is the best method on $13$ of $14$ columns in a controlled comparison under a
one-mask-per-query rule. The gains are largest on parts, up to $+14.5$ gIoU, and the model rejects more than
$90\%$ of the expressions that have no target. The same scores extend to reasoning queries and provide usable
masks for image editing.

\paragraph{Limitations and future work.}
\ours{} can only return a mask that the segmenter proposed, so its answers are bounded by the candidate pool.
Our analysis shows that this is not yet the limit: the best mask in the pool is far better than the one selected
($71.0$ against $49.4$ gIoU on PARTONOMY), so most of the remaining errors are errors of choice. Closing this gap is
the most direct way to improve the method. Reasoning queries need an extra rewriting step, and on object-level
reasoning \ours{} matches the box-based pipeline but does not exceed it. Finally, we studied requests with one
target; selecting several masks per request works with the same scores, and treating it as a first-class task is
left for future work.
